% !TeX root = ../../Book4.tex
% 纲要标题：### 16.4 群扩张谱序列
% 纲要标签：b4:sec:1504
% 纲要位置：计划纲要.md，第 3642 行。

\section{群扩张谱序列}
\label{b4:sec:1504}

群扩张的上同调可以先沿正规子群计算，再沿商群计算。商群在中间上同调上的作用与低阶传递映射，都是这一分步计算不可省略的数据。

% 纲要内容：
%
% * Lyndon–Hochschild–Serre 谱序列
% * 不变量函子的复合
% * 五项正合列与具体扩张的解释
% * 选择一个明确的有限循环群扩张写出五项正合列中的各项及传递映射；若使用连续群版本，必须转到附录 C 的连续模假设而非无标记替换
%
% ---
%

\subsection{Lyndon–Hochschild–Serre 谱序列}
\label{b4:subsec:1504:01}

设 \(1\to N\to G\xrightarrow{\pi}Q\to1\) 是离散群的扩张，\(M\) 为左 \(\mathbb ZG\)-模。研究 \(G\) 的上同调，可以先研究正规子群 \(N\)，再研究商群 \(Q\) 对所得上同调的作用。这由 Lyndon--Hochschild--Serre 谱序列\footnote{林登（Roger Conant Lyndon，1917-1988），美国数学家，研究群论与拓扑；塞尔（Jean-Pierre Serre，1926-），法国数学家，研究代数拓扑、代数几何与数论。霍赫希尔德已见本卷前注。}组织。第二步需要保留作用；即使已经算出底交换群 \(H^q(N,M)\)，也还不足以计算商群的上同调。

\begin{theorem}[Lyndon--Hochschild--Serre 谱序列]\label{b4:thm:lhs-spectral-sequence}
设 \(1\to N\to G\to Q\to1\) 是离散群扩张，\(M\) 为左 \(\mathbb ZG\)-模。则有自然的第一象限强收敛\term[lhs-puxulie]{Lyndon--Hochschild--Serre 谱序列}[Lyndon--Hochschild--Serre spectral sequence]
\[
 E_2^{p,q}=H^p(Q,H^q(N,M))\Longrightarrow H^{p+q}(G,M).
\]
其中 \(Q\) 在 \(H^q(N,M)\) 上的作用由 \(G\) 的共轭和系数作用共同诱导。其边缘同态分别为膨胀与限制所诱导的映射。
\end{theorem}

例如在一次上同调中，对余圈 \(c:N\to M\) 和 \(g\in G\)，作用由
\[
 (g\cdot c)(n)=g\cdot c(g^{-1}ng)
\]
给出：输入先共轭，输出再接受系数作用。\cref{b4:prop:normal-subgroup-invariants}已证明 \(N\) 自身的作用在上同调上平凡，故这才下降为 \(Q\) 作用；余链上的作用不必通过 \(Q\)。下一小节给出谱序列的证明。这里所有群都带离散结构；连续群上同调需要连续模及相应消解，不能直接把 \(G\) 换成拓扑群而保留全部假设。

\subsection{不变量函子的复合}
\label{b4:subsec:1504:02}

\begin{proof}[\cref{b4:thm:lhs-spectral-sequence}的证明]
在模范畴之间取
\[
 F:\mathbb ZG\text{-}\mathbf{Mod}\longrightarrow
 \mathbb ZQ\text{-}\mathbf{Mod},\quad F(M)=M^N,
 \qquad G_0=(-)^Q.
\]
\(N\) 正规保证 \(M^N\) 带 \(Q\) 作用，且 \(G_0F(M)=M^G\)。两个函子都加性左正合。

由\cref{b4:prop:normal-subgroup-invariants}，\((- )^N\) 是膨胀函子的右伴随；膨胀不改变底交换群的正合性，因而正合。用\cref{b4:prop:right-adjoint-injectives}，\(F\) 送内射 \(\mathbb ZG\)-模到内射 \(\mathbb ZQ\)-模，特别地其像对 \((- )^Q\) 零调。这一伴随核对的是复合函子定理所需的内射像条件。

取 \(G\)-内射消解 \(M\to I^\bullet\)。还要核对中间导出确实计算 \(N\) 的群上同调，这用的是另一套伴随：\(\mathbb ZG\) 作为右 \(\mathbb ZN\)-模自由，故诱导 \(\mathbb ZG\otimes_{\mathbb ZN}-\) 正合；其右伴随限制函子由同一命题保持内射对象。因此 \(I^\bullet\) 限制到 \(N\) 后仍是内射消解，得到底交换群的自然同构
\[
 R^qF(M)=H^q\bigl((I^\bullet)^N\bigr)\cong H^q(N,M).
\]

还须核对这一同构与 \(Q\) 作用相容。\cref{b4:prop:normal-subgroup-invariants}的证明用双复形
\[
 K^{i,j}=\operatorname{Hom}_{\mathbb ZN}(B_i(N),I^j)\qquad(i,j\ge0)
\]
给出两个增广拟同构
\[
 (I^\bullet)^N\longrightarrow\operatorname{Tot}(K)
 \longleftarrow C^\bullet(N,M).
\]
在 \(K\) 上，\(g\in G\) 对 bar 坐标作共轭，并作用于系数 \(I^j\)。bar 增广在共轭下不变，而 \(M\to I^0\) 是 \(G\)-线性的，所以两个比较映射都与 \(G\) 作用交换。它们诱导的上同调同构因此也是 \(G\) 等变的。左边 \(N\) 逐项平凡作用；右边的共轭余链作用在取上同调后才使 \(N\) 平凡作用，故两边都下降为 \(Q\) 作用。这就把底交换群的识别提升为 \(\mathbb ZQ\)-模的自然同构。

现在 Grothendieck 定理的全部假设已经核对。代入\cref{b4:thm:grothendieck-spectral-sequence}，第二页是所写群上同调，目标为不变量函子 \((- )^G\) 的导出，即 \(H^*(G,M)\)。底行的增广对应把 \(Q\)-余链沿 \(G\to Q\) 拉回并用 \(M^N\hookrightarrow M\) 作为系数映射，这正是膨胀；最左列的增广对应把 \(G\)-余链限制到 \(N\)，并由正规性落在 \(Q\)-不变类中。因此边缘同态是所述膨胀和限制。自然性与强收敛均由复合函子构造得到。
\end{proof}

这个证明没有要求群有限。有限性将在具体消解计算中使用，或者使某些平均算子可用；它不是谱序列存在的必要条件。标准叙述可参阅 Weibel\cite[Lemma 6.8.1, Lyndon/Hochschild-Serre Spectral Sequence 6.8.2]{Weibel1994HomologicalAlgebra}。

\subsection{五项正合列与群扩张}
\label{b4:subsec:1504:03}

由\cref{b4:prop:spectral-five-term}，每个离散群扩张及左 \(\mathbb ZG\)-模 \(M\) 给出
\begin{equation}\label{b4:eq:lhs-five-term}
\begin{aligned}
0&\longrightarrow H^1(Q,M^N)\xrightarrow{\mathrm{inf}}H^1(G,M)
\xrightarrow{\mathrm{res}}H^1(N,M)^Q\\
&\xrightarrow{\mathrm d_2^{0,1}}H^2(Q,M^N)
\xrightarrow{\mathrm{inf}}H^2(G,M).
\end{aligned}
\end{equation}
这里的微分 \(\mathrm d_2^{0,1}\) 称为\term[chuandi-yingshe]{传递映射}[transgression]。这个名称专指谱序列中的低阶微分，不是有限指数子群的 transfer（亦称余限制）。其含义是：一个 \(N\) 上的 \(Q\)-不变一次上同调类，若想成为 \(G\) 上一次类的限制，必须且只须其传递类消失。

\ProseParagraphBreak
若系数 \(A\) 的作用平凡，一次上同调就是群同态 \(G\to A\)。取满足 \(s(1)=1\) 的集合截面 \(s:Q\to G\)，记
\[
 c(x,y)=s(x)s(y)s(xy)^{-1}\in N.
\]
对 \(Q\)-不变同态 \(\phi:N\to A\)，定义 \(u(x,y)=\phi(c(x,y))\)。群乘法结合律给出
\[
 c(x,y)c(xy,z)=s(x)c(y,z)s(x)^{-1}c(x,yz).
\]
施加 \(\phi\)，利用它在共轭下不变，得到
\[
 u(x,y)+u(xy,z)=u(y,z)+u(x,yz),
\]
所以 \(u\) 是归一化 \(2\)-余圈。若 \(f:G\to A\) 延拓 \(\phi\)，令 \(b(x)=f(s(x))\)，则 \(u(x,y)=b(x)+b(y)-b(xy)=\delta b(x,y)\)。反过来，若 \(u=\delta b\)，则在唯一分解 \(g=ns(x)\) 下定义 \(f(g)=\phi(n)+b(x)\)，由共轭不变性及截面乘法公式可验证 \(f\) 为所需同态。因此 \([u]\) 检测延拓障碍。

\ProseParagraphBreak
具体符号还要与谱序列的代表元比较。取本节已经使用的 \(G\) 内射消解 \(A\xrightarrow{j}I^0\to I^1\to\cdots\)，记其微分为 \(\partial\)。\(I^0\) 限制到 \(N\) 后内射，故 \(H^1(N,I^0)=0\)。一次余圈 \(j\phi\) 因而是余边，可选 \(a\in I^0\) 使
\[
na-a=-j\phi(n)\quad(n\in N),\qquad \alpha=\partial a\in(I^1)^N.
\]
这里 \(n\alpha-\alpha=\partial(na-a)=0\)，且 \(\partial\alpha=0\)。在\cref{b4:prop:normal-subgroup-invariants}所用的 \(N\) 的 bar 消解与内射消解的比较双复形中，零次值 \(-a\) 的横、纵微分分别是 \(j\phi\) 与 \(-\alpha\)，所以
\[
\mathrm d(-a)=(j\phi,-\alpha),\qquad
(j\phi,0)-\mathrm d(-a)=(0,\alpha).
\]
因此 \(\alpha\) 对应的是 \([\phi]\)，而不是它的负类。

\ProseParagraphBreak
用 \(Q\) 的 bar 消解计算第二步，双复形可写为
\[
D^{q,p}=C^p\bigl(Q,(I^q)^N\bigr),\qquad
\mathrm d_h=\partial,\qquad \mathrm d_v=(-1)^q\delta_Q.
\]
第一指标 \(q\) 是原消解次数，第二指标 \(p\) 是商群余链次数；谱序列按第二指标滤过，第二页则写作 \(E_2^{p,q}\)。各 \((I^q)^N\) 都是 \(Q\) 内射模，bar 计算与内射计算经消解比较给出前面的复合函子谱序列。以上总化符号与\cref{b4:def:hypercohomology}相同。令 \(\beta_x=s(x)a-a\)。由 \(\phi\) 的共轭不变性，
\[
n\beta_x-\beta_x
=s(x)\bigl((s(x)^{-1}ns(x))a-a\bigr)-(na-a)=0,
\]
所以 \(\beta\in C^1(Q,(I^0)^N)\)，并且 \(\delta_Q\alpha=\partial\beta\)。\(\alpha\) 位于 \((q,p)=(1,0)\)，其商群方向的微分带负号；\(\beta\) 位于 \((0,1)\)，该方向的微分带正号。修正代表元后有
\[
\mathrm d(\alpha+\beta)
=-\delta_Q\alpha+\partial\beta+\delta_Q\beta
=\delta_Q\beta.
\]
由截面的乘法公式，再用 \(A\) 的平凡作用及 \(\phi\) 的共轭不变性，得到
\[
\begin{aligned}
(\delta_Q\beta)(x,y)
&=\bigl(s(x)s(y)-s(xy)\bigr)a\\
&=s(xy)\bigl((s(xy)^{-1}c(x,y)s(xy))a-a\bigr)\\
&=-j\phi(c(x,y))=-j u(x,y).
\end{aligned}
\]
它落在 \(\ker\partial=j(A)\)，正是第二页目标的系数。把每个 \(D^{q,p}\) 分量的元素乘以 \((-1)^{pq}\)，就得到从当前总复形到另一总复形的同构：商群方向的微分 \((-1)^q\delta_Q\) 变为通常的 \(\delta_Q\)，消解方向的微分则变为 \((-1)^p\partial\)。这个同构保持按 \(p\) 的滤过；在第二页的源 \((p,q)=(0,1)\) 和靶 \((2,0)\) 处，它的系数都为一。因此由\cref{b4:thm:filtered-spectral-sequence}的代表元规则，得到
\[
\mathrm d_2^{0,1}[\phi]=-[u].
\]
这里使用了平凡系数，任意系数时不能省去交叉同态的作用项。下一小节在有限循环群中计算这项障碍。低阶正合列的教材叙述见 Weibel\cite[Low Degree Terms 6.8.3]{Weibel1994HomologicalAlgebra}。

\subsection{有限循环群扩张的传递映射}
\label{b4:subsec:1504:04}

取
\[
 1\longrightarrow N=\langle t^2\rangle\longrightarrow
 G=\langle t\mid t^4=1\rangle\longrightarrow Q=C_2\longrightarrow1,
 \qquad M=\mathbb F_2
\]
并令系数作用平凡。所有共轭作用也平凡。由\cref{b4:thm:cyclic-group-cohomology}，五项列中五个非零候选项均为 \(\mathbb F_2\)：
\[
 0\longrightarrow\mathbb F_2\xrightarrow{\mathrm{inf}}
 \mathbb F_2\xrightarrow{\mathrm{res}}\mathbb F_2
 \xrightarrow{\mathrm d_2^{0,1}}\mathbb F_2
 \xrightarrow{\mathrm{inf}}\mathbb F_2.
\]
第一次膨胀把 \(Q\) 的生成元映到一的同态拉回到 \(t\mapsto1\)，所以为恒等同构。限制却为零，因为任意 \(f:G\to\mathbb F_2\) 满足 \(f(t^2)=2f(t)=0\)。正合性已迫使 \(\mathrm d_2^{0,1}\) 为同构，最后一次膨胀为零。

\ProseParagraphBreak
还可以写出传递类的余圈。用加法记 \(Q=\{0,1\}\)，取 \(s(0)=1\)、\(s(1)=t\)。令 \(\phi:N\to\mathbb F_2\) 由 \(\phi(t^2)=1\) 定义，则
\[
 c(x,y)=s(x)s(y)s(x+y)^{-1}\in N,
 \qquad u(x,y)=\phi(c(x,y))
\]
仅在 \((x,y)=(1,1)\) 取一。直接代入余圈等式
\(u(y,z)-u(x+y,z)+u(x,y+z)-u(x,y)=0\)，或者用截面乘法的结合律，即可验证 \(u\) 是\(2\)-余圈。它不是余边：若 \(u=\delta v\)，则 \(0=u(0,y)=\delta v(0,y)=v(0)\)，而 \(\delta v(1,1)=v(1)-v(0)+v(1)=0\)，与 \(u(1,1)=1\) 矛盾。\(H^2(C_2,\mathbb F_2)\) 只有一个非零元，因此已证明非零的传递映射 \(\mathrm d_2^{0,1}\) 把 \([\phi]\) 送到 \([u]\)。本例特征为二，\(-[u]=[u]\)，与上述符号公式一致。

\ProseParagraphBreak
为直接核对最后一次膨胀为零，对 \(0\le i\le3\) 定义一余链 \(b(t^i)=\lfloor i/2\rfloor\pmod2\)。把 \(i=2a+r\)、\(j=2a'+r'\)（\(r,r'\in\{0,1\}\)）代入，得到
\[
 \delta b(t^i,t^j)=b(t^j)-b(t^{i+j})+b(t^i)=rr'
 =u(\pi(t^i),\pi(t^j)).
\]
指数取模四不影响右边的模二计算。因此膨胀后的\(2\)-余圈确实成为余边。与分裂扩张 \(C_2\to C_2\times C_2\to C_2\) 比较，后一情形每个 \(N\to\mathbb F_2\) 都通过第一因子延拓，限制为满，故传递映射为零。这个差别说明低阶微分记录了群扩张本身的信息，而不只由核群与商群的同构类型决定。

\ProseParagraphBreak
要使负号可见，取 \(G=C_9=\langle t\mid t^9=1\rangle\) 和扩张 \(1\to\langle t^3\rangle\to G\to C_3\to1\)，系数为平凡的 \(\mathbb F_3\)。用 \(0,1,2\) 表示商群元素，取 \(s(i)=t^i\)、\(\phi(t^3)=1\)。截面因子集给出进位余圈
\[
u(i,j)=\left\lfloor\frac{i+j}{3}\right\rfloor\pmod3,
\qquad \mathrm d_2^{0,1}[\phi]=-[u].
\]
这两种符号确实代表不同的类。对任意一余链 \(v:C_3\to\mathbb F_3\)，
\[
\sum_{i=0}^{2}(\delta v)(i,1)
=\sum_{i=0}^{2}\bigl(v(1)-v(i+1)+v(i)\bigr)
=3v(1)=0,
\]
其中 \(i+1\) 在 \(C_3\) 中计算。而 \(\sum_i u(i,1)=1\)、\(\sum_i(-u(i,1))=2\)，所以 \([u]\ne-[u]\)。若以 \([\phi]\) 和 \([u]\) 分别作为两端一维空间的基，本书的传递映射就是乘 \(-1=2\) 的同构。
